\begin{proof}
    Let $\{ b_j \}_{j=1}^n$ denote the orthonormal eigenbasis associated with $\{ \lambda_j \}_{j=1}^n$.
    Then let $\alpha \in \mathbb{R}$ and $\beta \in \mathbb{R}$.
    First, equations will be found for $( U^\alpha )^\beta$ and $U^{\alpha\beta}$.
    Observe that the following equation holds.
    \begin{equation*}
        U^\alpha
        =
        \exp( \alpha\Log( U ) )
        =
        \exp\left( \alpha \sum_{j=1}^n i\lambda_j\pi \ket{b_j} \bra{b_j} \right)
        =
        \exp\left( \sum_{j=1}^n i\alpha\lambda_j\pi \ket{b_j} \bra{b_j} \right)
        =
        \sum_{j=1}^n e^{i\alpha\lambda_j\pi} \ket{b_j} \bra{b_j}
    \end{equation*}
    This can then be used to compute $\Log( U^\alpha )$ where $\zeta_j = \Log\left( e^{\alpha\lambda_j} \right)$.
    \begin{equation*}
        \Log( U^\alpha )
        =
        \Log\left( \sum_{j=1}^n e^{i\alpha\lambda_j\pi} \ket{b_j} \bra{b_j} \right)
        =
        \sum_{j=1}^n \Log\left( e^{i\alpha\lambda_j\pi} \right) \ket{b_j} \bra{b_j}
        =
        \sum_{j=1}^n \zeta_j \ket{b_j} \bra{b_j}
    \end{equation*}
    This can then be used to compute $( U^\alpha )^\beta$.
    \begin{equation*}
        ( U^\alpha )^\beta
        =
        \exp( \beta\Log( U^\alpha ) )
        =
        \exp\left( \beta \sum_{j=1}^n \zeta_j \ket{b_j} \bra{b_j} \right)
        =
        \exp\left( \sum_{j=1}^n \beta \zeta_j \ket{b_j} \bra{b_j} \right)
        =
        \sum_{j=1}^n e^{\beta\zeta_j} \ket{b_j} \bra{b_j}
    \end{equation*}
    Likewise, $U^{\alpha\beta} = \sum_{j=1}^n e^{i\alpha\beta\lambda_j\pi} \ket{b_j} \bra{b_j}$.
    Next an explicit equation will be calculated for each $\zeta_j$ in terms of $\alpha$ and $\lambda_j$.
    Let $j \in \{ 1, 2, \ldots, n \}$.
    Since $\zeta_j = \Log( \exp( i\alpha\lambda_j\pi ) )$, then there exists a $k \in \mathbb{Z}$ such that $i\alpha\lambda_j\pi = \zeta_j + k(i2\pi)$.
    Then the following sequence of statements hold.
    \begin{align*}
        \zeta_j
        &\in
        i( -\pi, \pi ]
        &
        ( \text{since } \zeta_j = \Log( \exp( \alpha\lambda_j ) ) )
        \\
        i\alpha\lambda_j\pi
        &\in
        i( -\pi + k(2\pi), \pi + k(2\pi) ]
        &
        ( \text{since } \alpha\lambda_j = \zeta_j + k(i2\pi) )
        \\
        \alpha\lambda_j / 2
        &\in
        ( k - 1/2, k + 1/2 ]
        \\
        \alpha\lambda_j / 2 - 1/2
        &\in
        ( k - 1, k ]
    \end{align*}
    Since $k - 1 < ( \alpha\lambda_j - 1 ) / 2 \le k$, then $\ceil*{( \alpha \lambda_j - 1 ) / 2} = k$.
    Since $j$ was arbitrary, then it follows that $\alpha\lambda_j = \zeta_j + \ceil*{( \alpha \lambda_j - 1 ) / 2} (i2\pi)$ for each $j \in \{ 1, 2, \ldots, n \}$.
    Next, it must be shown that $( U^\alpha )^\beta = U^{\alpha\beta}$ if and only if for each $j \in \{ 1, 2, \ldots, n \}$, either $\alpha \lambda_j \in ( -1, 1 ]$ or $\beta \in \mathbb{Q}$ with reduced denominator dividing $\ceil*{( \alpha \lambda_j - 1 ) / 2}$.
    \begin{itemize}
    \item Assume that $U^{\alpha\beta} = ( U^\alpha )^\beta$.
          Let $j \in \{ 1, 2, \ldots, n \}$.
          Since $\{ b_j \}_{j=1}^n$ is an orthonormal basis, then $e^{\beta\zeta_j} = e^{i\alpha\beta\lambda_j\pi}$ by equating components in $( U^\alpha )^\beta = U^{\alpha\beta}$.
          Then $\beta\zeta_j - i\alpha\beta\lambda_j\pi$ is a multiple of $i2\pi$.
          Then there exists some $m \in \mathbb{Z}$ such that $m (i2\pi) = \beta(i\alpha\lambda_j\pi - \zeta_j) = \beta k (i2\pi)$ where $k = \ceil*{( \alpha \lambda_j - 1 ) / 2}$.
          In other words, $\beta k \in \mathbb{Z}$.
          If $k = 0$, then clearly $\beta k \in \mathbb{Z}$.
          This happens precisely when $\zeta_j = i\alpha\lambda_j\pi$, which implies that $\alpha\lambda_j \in (-1,1]$.
          Assume instead that $k \ne 0$.
          Then $\beta \in \mathbb{Q}$, otherwise $\beta k$ is irrational.
          Then there exists some $d \in \mathbb{Z}$ and $q \in \mathbb{N}$ such that $\gcd( d, q ) = 1$ and $\beta = d / q$.
          Then $dk / q \in \mathbb{Z}$.
          Then $q$ divides $dk$.
          Since $\gcd( q, d ) = q$, then $q$ divides $k$.
          This means that $\lambda_j$ is $( \alpha, \beta )$-admissible.
          Since $j$ was arbitrary, then $\lambda_j$ is $( \alpha, \beta )$-admissible for each $j \in \{ 1, 2, \ldots, n \}$.
    \item Assume that for each $j \in \{ 1, 2, \ldots, n \}$, $\lambda_j$ is $( \alpha, \beta )$-admissible.
          Let $j \in \{ 1, 2, \ldots, n \}$.
          There are two cases to consider.
          If $\alpha \lambda_j \in (-1,1]$, then $\zeta_j = i\alpha \lambda_j\pi$, and consequently $\beta\zeta_j = i\alpha\beta\lambda_j\pi = i\alpha\beta\pi + 0(i2\pi)$.
          Assume instead that $\beta \in \mathbb{Q}$ with reduced denominator dividing $k = \ceil*{( \alpha \lambda_j - 1 ) / 2}$.
          Then there exists $q \in \mathbb{Z}$ and $d \in \mathbb{N}$ such that $\beta = q / d$, $\gcd( q, d ) = 1$ and $d$ divides $k$.
          It follows that $\beta\zeta_j = \beta( i\alpha\lambda_j\pi + k ) = i\alpha\beta\lambda_j\pi + q k (i2\pi) / d$.
          Since $d$ divides $k$, then there exists a $m \in \mathbb{Z}$ such that $\beta\zeta_j = i\alpha\beta\lambda_j\pi + m(i2\pi)$.
          In either case, there exists an $m \in \mathbb{N}$ such that $\beta\zeta_j = i\alpha\beta\lambda_j\pi + m(i2\pi)$.
          Then $\exp( \beta\zeta_j ) = \exp( i\alpha\beta\lambda_j\pi )$.
          Since $j$ was arbitrary, then $\exp( \beta\zeta_j ) = \exp( i\alpha\beta\lambda_j\pi )$ for each $j \in \{ 1, 2, \ldots, n \}$.
          Then the following equation holds.
          \begin{equation*}
              ( U^\alpha )^\beta
              =
              \sum_{j=1}^n e^{\beta\zeta_j} \ket{b_j} \bra{b_j}
              =
              \sum_{j=1}^n e^{i\alpha\beta\lambda_j\pi} \ket{b_j} \bra{b_j}
              =
              U^{\alpha\beta}
          \end{equation*}
    \end{itemize}
    Then $( U^\alpha )^\beta = U^{\alpha\beta}$ if and only if for each $j \in \{ 1, 2, \ldots, n \}$, either $\alpha \lambda_j \in ( -1, 1 ]$ or $\beta \in \mathbb{Q}$ with reduced denominator dividing $\ceil*{( \alpha \lambda_j - 1 ) / 2}$.
    It remains to be shown that $\mathbb{Z}$ and $( -1, - 1]$ are maximal submonoids of $\mathbb{R}$ for which for which $( U, \alpha ) \mapsto U^\alpha$ is a monoid action on $\Mor( \CatUnitary )$.
    \begin{itemize}
    \item Let $\alpha \in \mathbb{Z}$, $\beta \in \mathbb{Z}$, and $U \in \Mor( \CatUnitary )$.
          It must be shown that $( U^\alpha )^\beta = U^{\alpha\beta}$.
          To this end, let $\{ \lambda_j \}_{j=1}^n$ denote the eigenvalues of $\Log( U ) / (i\pi)$.
          Let $j \in \{ 1, 2, \ldots, n \}$.
          Assume that $\alpha\lambda_j \not \in ( -1, 1 ]$.
          Since $\beta \in \mathbb{Z}$, then its reduced denominator is $1$.
          Then clearly the reduced denominator of $\beta$ divides $\ceil*{( \alpha \lambda_j - 1 ) / 2}$.
          Then either $\alpha\lambda_j \in ( -1, 1 ]$ or $\beta \in \mathbb{Q}$ with reduced denominator dividing $\ceil{( \alpha \lambda_j - 1) / 2 }$.
          Since $j$ was arbitrary, then $( U^\alpha )^\beta = U^{\alpha\beta}$.
          Since $\alpha$, $\beta$, and $U$ were arbitrary, then $( U, \alpha ) \mapsto U^\alpha$ is a monoid action on $\Mor( \CatUnitary )$ with respect to $\mathbb{Z}$.
          It remains to be shown that $\mathbb{Z}$ is maximal.
          Assume that $M$ is a submonoid of $\mathbb{R}$ for which $( U, \alpha ) \mapsto U^\alpha$ is a monoid action on $\Mor( \CatUnitary )$ and $\mathbb{Z}$ is a submonoid of $M$.
          Let $\beta \in M$.
          Then $( Z^2 )^\beta = Z^{2\beta}$ by definition of $M$.
          Since $\lambda = 1$ is an eigenvalue of $\Log( Z ) / (i\pi)$ and $2\lambda \not \in (-1,1]$, then it follows by the first result of this proof that $\beta \in \mathbb{Q}$.
          Then there exists some $q \in \mathbb{Z}$ and $d \in \mathbb{N}$ such that $\beta = q / d$ and $\gcd( q, d ) = 1$.
          Since $2d + 3 \in \mathbb{Z} \subseteq M$, then $( Z^{2d + 3} )^\beta = Z^{\beta(2d+3)}$ by definition of $M$.
          Since $\lambda = 1$ is an eigenvalue of $\Log( Z ) / (i\pi)$ and $(2d+3) \not \in (-1,1]$, then by the first result of this proof $d$ divides $\ceil{((2d+3)-1)/2} = d + 1$.
          This is only possible when $d = 1$.
          Then $\beta = q / d = q \in \mathbb{Z}$.
          Since $\beta$ was arbitrary, then $M \subseteq \mathbb{Z}$.
          In conclusion, $\mathbb{Z}$ is maximal.
    \item Let $\alpha \in ( -1, 1 ]$, $\beta \in ( -1, 1 ]$, and $U \in \Mor( \CatUnitary )$.
          It must be shown that $( U^\alpha )^\beta = U^{\alpha\beta}$.
          To this end, let $\{ \lambda_j \}_{j=1}^n$ denote the eigenvalues of $\Log( U ) / (i\pi)$.
          Since $\alpha \in ( -1, 1 ]$ and $i\lambda_j\pi \in ( -\pi, \pi ]$ for each $j \in \{ 1, 2, \ldots, n \}$, then $\alpha\lambda_j \in ( -1, 1 ]$ for each $j \in \{ 1, 2, \ldots, n \}$.
          Then by the first result of this proof, $( U^\alpha )^\beta = U^{\alpha\beta}$ regardless of the value of $\beta$.
          It remains to be shown that $( -1, 1 ]$ is maximal.
          Assume that $M$ is a submonoid of $\mathbb{R}$ for which $( U, \alpha ) \mapsto U^\alpha$ is a monoid action on $\Mor( \CatUnitary )$ and $( -1, 1 ]$ is a submonoid of $M$.
          Let $\alpha \in \mathbb{R}$.
          There are three cases to consider.
          \begin{enumerate}
          \item If $\alpha \in ( -1, 1 ]$, then $\alpha \in M$ by definition.
          \item If $\alpha = -1$, then $( Z^{\alpha} )^{1/2} = ( Z^{-1} )^{1/2} = Z^{1/2} = S \ne S^\dagger = Z^{-1/2} = Z^{\alpha/2}$.
                Since $1/2 \in M$, then $\alpha \not \in M$, since $M$ acts on $\Mor( \CatUnitary )$ by exponentiation.
          \item Assume that $\alpha \not \in [ -1, 1 ]$.
                Then $\alpha^2 > 1$ and there exists some $m \in \mathbb{N}$ such that $\alpha^{2m} > 2$.
                Let $\hat{\alpha} = \alpha^{2m}$.
                Then $( \hat{\alpha} - 1 ) / 2 > 1 / 2$.
                Then $\ceil*{( \hat{\alpha} / 2} \ge 1$.
                Let $k = \ceil*{( \hat{\alpha} - 1 ) / 2}$ and $\beta = 1 / ( k + 1 )$.
                Since $\lambda = 1$ is an eigenvalue of $\Log( Z ) / (i\pi)$ and $\hat{\alpha}\lambda = \hat{\alpha} \not \in ( -1, 1 ]$, then it follows from the first result of this proof that $( U{\hat{\alpha}} )^\beta = U^{\hat{\alpha}\beta}$ if and only if $k + 1$ divides $k$.
                Since $k \ne 0$ and $k + 1 > k$, then clearly $k + 1$ does not divide $k$.
                This means that $( U^{\hat{\alpha}} )^\beta \ne U^{\hat{\alpha}\beta}$.
                Since $M$ acts on $\Mor( \CatUnitary )$ by exponentiation and $\beta \in M$, then $\hat{\alpha} \not \in M$.
                Then $\alpha \not \in M$, since otherwise $\hat{\alpha} = \alpha^{2m} \in M$.
          \end{enumerate}
          This means that if $\alpha \in M$, then $\alpha \in ( -1, 1 ]$.
          In other words, $M \cap \mathbb{R} = ( -1, 1 ]$.
          However, $M$ was a submonoid of $\mathbb{R}$, so $M = ( -1, 1 ]$.
          In conclusion, $( -1, 1 ]$ is maximal.
    \end{itemize}
    Then $\mathbb{Z}$ and $( -1, 1 ]$ are maximal submonoids of $\mathbb{R}$ for which for which $( U, \alpha ) \mapsto U^\alpha$ is a monoid action on $\Mor( \CatUnitary )$.
\end{proof}